\section{Introduction}\label{sec:intro}

Inorganic scale such as calcium carbonate and barium sulphate precipitates when produced water becomes
supersaturated, for example when pressure falls and CO$_2$ leaves solution or when injected and formation
waters mix \citep{kamal2018,olajire2015}. Inside a production choke the deposit narrows the throat, raises
the upstream pressure and forces the operator to choke the well back; left untreated it ends in an
unplanned cleaning or acid job and deferred production. Because scale grows gradually, over hours to
weeks, a forecast of \emph{when} the restriction will become critical would allow the intervention to be
planned rather than improvised. In practice the restriction is still tracked manually from trend charts,
which makes choke scaling a natural target for the data-driven predictive maintenance that is now common in
manufacturing and process industries.

Public data for this problem became available with the 3W dataset released by Petrobras
\citep{vargas2019,vargas2026}, which contains labelled multivariate sensor series of real, simulated and
hand-drawn undesirable events in offshore wells, among them scaling in the production choke. The
machine-learning work on 3W has, however, treated the task almost exclusively as detection: a model
classifies whether an event is under way, usually after it has started \citep{aranha2024,lopes2025,lopes2026}.
Detection answers whether the well is scaling now, not how long it will take until the restriction becomes
critical. Evaluations also commonly split windows or instances at random, so that windows of the same well
appear in training and test sets; a recent community proposal for 3W shows that the macro-F1 score falls
from 0.99 under a random split to 0.30 when whole wells are held out \citep{issue231}. Outside 3W, the
data-driven studies that estimate the magnitude of scale rely on simulated well histories or on water
chemistry rather than on production-sensor streams of real wells
\citep{cheng2022,askar2025,su2024,molina2025}, and few studies attach calibrated uncertainty to their
outputs, although an alarm that is right on average but uncalibrated on a new well is of little use to an
operator. To the best of our knowledge, no published study forecasts the time to a critical level of choke
scaling from production sensors, evaluates such a forecast on wells never seen in training, or reports its
uncertainty.

This paper addresses that gap. Choke-scale forecasting is cast as discrete-time survival analysis, in which
every window of sensor data receives the probability that the restriction reaches a critical level within
horizons of 1 to 72~h and instances that never reach it are treated as right-censored. Because the
pressure downstream of the choke is missing in all real scaling instances of 3W, the restriction is measured
by an index normalised to each instance's own early baseline, which serves both as label and as causal input
and makes wells of very different pressure comparable. On this formulation we build PI-PatchSurv, a
channel-independent patch Transformer whose sequence representation is fused through a learned gate with a
branch of physics features and decoded by a survival head that yields a monotone risk curve; a monotonicity
loss, a deep ensemble and split-conformal intervals complete the model. The approach is evaluated against
eleven alternatives, from a physics threshold rule to re-implementations of the architectures of the base
studies, on leave-one-well-out folds in which every threshold and hyper-parameter is chosen on separate
calibration wells.

The novelty of this research lies in three aspects. First, it provides, as far as we are aware, the first
time-to-event formulation of an undesirable event in the 3W dataset, together with a self-referenced
restriction proxy that makes the formulation usable when the downstream sensors are missing. Second, it
combines learned sequence features, physical indicators and calibrated uncertainty in one survival model
for well operations. Third, it reports an unseen-well benchmark of twelve models with fixed and tuned
hyper-parameters, event-level alarm metrics, robustness tests and a dataset-version generalisation test, and
documents two pitfalls of 3W version 2.0.0 that silently change results. The results show that
time-to-restriction forecasting on unseen wells is feasible but that its lead time is bounded by the data,
and that simple recurrent forecasters rank events as well as the proposed model, whose advantages lie in
uncertainty quality and robustness.

The remainder of this paper is organised as follows. Section~\ref{sec:related} reviews related work, and
Section~\ref{sec:method} presents the proposed method. Section~\ref{sec:setup} describes the experimental
setup, Section~\ref{sec:results} reports the results, Section~\ref{sec:comparison} compares them with the
base studies, and Section~\ref{sec:ablation} presents the ablation study. Section~\ref{sec:discussion}
discusses implications and limitations, and Section~\ref{sec:conclusion} concludes the paper.
